% ==============================================================================
% APPENDIX A: XML Co-Pilot Command Schema
% ==============================================================================

\chapter{XML Co-Pilot Command Schema}
\label{app:xml_commands}

\noindent
This appendix lists the complete set of XML command schemas used by the chatbot co-pilot thread to communicate edit actions back to the PySide6 symbolic editor canvas view. To enable real-time, non-blocking command execution, the co-pilot orchestrator streams natural language responses from the LLM, extracts structured XML blocks using regular expressions, and dispatches them sequentially to the main GUI thread's command queue.

\section{Command Parsing and Dispatching Pipeline}
\noindent
When the orchestrator worker thread receives a streamed response from the LLM, it extracts the command packet using a regular expression that targets blocks enclosed in \texttt{<CMD>} tags. The extracted block is parsed using Python's standard \texttt{xml.etree.ElementTree} library. 

\noindent
The parser extracts the action type from the \texttt{<action>} node and routes the payload to the layout tab controller (\texttt{layout\_tab.py}), which maps the XML values to the command stack. The command stack instantiates a corresponding undo-aware command object (inheriting from \texttt{QUndoCommand}) and pushes it onto the application's active \texttt{QUndoStack}, triggering a canvas redraw and updating the head-less KLayout physical view.

\section{Symbolic Movement Command (MOVE)}
\noindent
The \texttt{MOVE} command shifts a specified device instance horizontally or vertically by a relative number of symbolic grid slot units.
\begin{thesiscode}{Transistor Move Schema}{xml}
<CMD>
  <action>MOVE</action>
  <instance_id>MM1</instance_id>
  <x_slot>2</x_slot>
  <y_row>1</y_row>
</CMD>
\end{thesiscode}

\section{Device Instance Swapping Command (SWAP)}
\noindent
The \texttt{SWAP} command exchanges the symbolic positions of two transistor instances on the placement canvas, which is frequently used to optimize matching orientations.
\begin{thesiscode}{Transistor Swap Schema}{xml}
<CMD>
  <action>SWAP</action>
  <instance_a>MM1</instance_a>
  <instance_b>MM2</instance_b>
</CMD>
\end{thesiscode}

\section{Active Diffusion Abutment Command (ABUT)}
\noindent
The \texttt{ABUT} command groups two adjacent devices together, aligning their coordinates and setting PDK parameter flags to merge their active diffusions.
\begin{thesiscode}{Active Region Abutment Schema}{xml}
<CMD>
  <action>ABUT</action>
  <instance_a>MM1</instance_a>
  <instance_b>MM2</instance_b>
  <abut_type>shared</abut_type>
</CMD>
\end{thesiscode}

\section[Multi-Device Alignment (ALIGN)]{Instance Multi-Device Alignment\\Command (ALIGN)}
\noindent
The \texttt{ALIGN} command aligns a list of device instances along a specified geometric axis (horizontal or vertical) relative to a reference coordinate.
\begin{thesiscode}{Instance Align Schema}{xml}
<CMD>
  <action>ALIGN</action>
  <instances>
    <id>MM1</id>
    <id>MM2</id>
    <id>MM3</id>
  </instances>
  <axis>vertical</axis>
  <alignment>center</alignment>
</CMD>
\end{thesiscode}

\section{Transistor Orientation Mirroring Command (ORIENT)}
\noindent
The \texttt{ORIENT} command mirrors the gate poly fingers of an instance (setting rotation and source-drain mirroring orientations like \texttt{R0}, \texttt{MY}, or \texttt{MX}).
\begin{thesiscode}{Transistor Orientation Schema}{xml}
<CMD>
  <action>ORIENT</action>
  <instance_id>MM1</instance_id>
  <orientation>MY</orientation>
</CMD>
\end{thesiscode}

\section[Symmetry Enforcement (SYMMETRY)]{Symmetry Constraint Enforcement\\Command (SYMMETRY)}
\noindent
The \texttt{SYMMETRY} command binds a list of matched pairs symmetrically about a specific midpoint axis, triggering reflection checks in the symmetry enforcer.
\begin{thesiscode}{Symmetry Enforcement Schema}{xml}
<CMD>
  <action>SYMMETRY</action>
  <group_id>diff_pair_1</group_id>
  <midpoint>0.294</midpoint>
</CMD>
\end{thesiscode}

\section{DRC Auto-Heal Correction Command (DRC\_HEAL)}
\noindent
The \texttt{DRC\_HEAL} command applies absolute spacing correction shifts generated by the DRC Critic node to resolve clearance or overlapping violations.
\begin{thesiscode}{DRC Auto-Heal Schema}{xml}
<CMD>
  <action>DRC_HEAL</action>
  <instances>
    <instance>
      <id>MM1</id>
      <x_shift>0.140</x_shift>
      <y_shift>0.000</y_shift>
    </instance>
    <instance>
      <id>MM2</id>
      <x_shift>-0.140</x_shift>
      <y_shift>0.000</y_shift>
    </instance>
  </instances>
</CMD>
\end{thesiscode}

\section{Physical Compilation and Export Command (EXPORT)}
\noindent
The \texttt{EXPORT} command triggers physical layout stream generation, compiling the current placement into an OASIS/GDSII file or writing updates to OpenAccess.
\begin{thesiscode}{Physical Export Schema}{xml}
<CMD>
  <action>EXPORT</action>
  <format>OASIS</format>
  <cell_name>comparator_latch</cell_name>
</CMD>
\end{thesiscode}

\section{Command Stack Undo Command (UNDO)}
\noindent
The \texttt{UNDO} command rolls back the last command pushed onto the PySide6 command stack, restoring the previous coordinates and orientations from memory.
\begin{thesiscode}{Command Stack Undo Schema}{xml}
<CMD>
  <action>UNDO</action>
</CMD>
\end{thesiscode}
